% SPDX-License-Identifier: Apache-2.0
\section{Building This Document}
\label{sec:building}

This article is a build artifact. It is compiled from LaTeX by Bazel, and
published by the repository's release workflow alongside every other
document.

\subsection{A hermetic LaTeX toolchain}

The build takes nothing from the machine it runs on. The rules resolve a
toolchain type rather than a binary, and two implementations exist: a system
toolchain that wraps the host's installed tools, and a hermetic one built
from pinned archives. The root module registers the hermetic toolchain,
which therefore wins, and the build needs nothing installed.

Three archives make it up. A TeX distribution published as a small
self-contained per-platform tarball, whose binaries locate their own tree
with no install step. A PDF utility that backs the page-count and
concatenation contract. Ghostscript, used only for the outline of a combined
PDF. Each is pinned by URL and by SHA-256.

That distribution includes neither IEEEtran nor the drawing packages, so the
missing ones are added as content-addressed package archives from a dated
snapshot of the TeX Live network repository. The alternative mechanism
resolves dependencies for you and talks to a live mirror, which stops the
fetch being reproducible, so every dependency is listed explicitly instead.

The snapshot date matches the vintage of the pinned distribution. This is
not decoration. TeX Live is a rolling release with one live repository, and
the moment upstream updates the package manager's own infrastructure, an
older package manager pointed at a live mirror refuses to run. A
distribution pinned by hash plus a repository that moves is a fetch that
breaks on a schedule.

\subsection{Why not Markdown}

An earlier arrangement rendered the documents from Markdown through pandoc,
with the whole tool set provisioned from a Debian package snapshot. It works
and it is hermetic, and it does not run on the project's build machine,
which is why this article is LaTeX.

\subsection{Why the figures are text and not TikZ}

Every figure here is a framed listing rather than a drawing, and the reason
is the toolchain rather than a preference. TikZ is the right tool for a
box-and-arrow diagram in LaTeX. It is not available: the pinned TeX
distribution includes no \code{pgf}, and the package cannot be fetched
reproducibly for the reason Section~\ref{sec:building} already gives. Adding
it means vendoring 6.4\,MB across 260 files.

Text figures do buy one thing. A TikZ picture can be wrong in a way that
produces no warning and no error: the document compiles clean, every
reference resolves, no box is overfull, and an arrowhead is still buried in
a box border. Only rendering the page and looking at it finds that. A text
figure is already what it will look like.

That is a reason to be comfortable with the current figures, not a reason to
prefer them. When \code{pgf} is vendored, the diagrams should become TikZ.

\subsection{Release}

The release workflow builds every document and publishes the results. A
nightly schedule overwrites a rolling release, so one link always points at
the current documents. A manual run instead cuts a release stamped with the
date and time, which nothing later overwrites.

\section{Conclusion}
\label{sec:conclusion}

Two languages written independently for one project turned out to be two
halves of one language. LHDL specifies how a system is assembled,
configured and scheduled, and cannot express a protocol body. TxHDL
specifies how a module talks a protocol, and cannot express a system. The
merge keeps almost all of both.

Only four questions needed a decision about meaning. The model of time is
the one that shapes the language, and the answer is that both models are
right about different things: a number that comes from a protocol on a wire
is a specification, and a number that comes from how fast the fabric
happens to be is not.

Everything else was spelling, and following Rust settled it while removing
six defects the sources had. The remaining defects are the ones nothing
checked, and the merged specification is being written so that a machine
checks them.
